\documentclass[12pt,a4paper]{report}

% --- Paquetes de Idioma y Codificación ---
\usepackage[utf8]{inputenc}
\usepackage[T1]{fontenc}
\usepackage[spanish, es-tabla]{babel}

% --- Diseño de Página ---
\usepackage{geometry}
\geometry{top=2.5cm, bottom=2.5cm, left=3cm, right=3cm}

% --- Matemáticas y Símbolos ---
\usepackage{amsmath, amssymb, amsthm}

% --- Colores y Cajas ---
\usepackage[dvipsnames]{xcolor}
\usepackage[most]{tcolorbox}

% Caption

\usepackage{caption}

% Gráficos

\usepackage{graphicx}

% Espaciadora

\newcommand{\esp}{\vspace{0.5cm}}

% promedio

\newcommand{\prom}[1]{\left\langle #1 \right\rangle}

% ket

\newcommand{\ket}[1]{| #1 \rangle}

% bra

\newcommand{\bra}[1]{\langle #1 |}

% braket

\newcommand{\braket}[2]{\langle #1 | #2 \rangle}

% Bibliografía

\usepackage[style=apa, backend=biber]{biblatex} % Estilo APA
\addbibresource{biblio.bib} % Nombre de tu archivo .bib

% Configuración de cajas personalizadas
\newtcolorbox{importante}{
	colback=red!5!white,
	colframe=red!75!black,
	title=Importante,
	fonttitle=\bfseries
}

\newtcolorbox{nota}{
	colback=blue!5!white,
	colframe=blue!75!black,
	title=Nota,
	fonttitle=\bfseries
}

% Definición de la caja de bibliografía
\newtcolorbox{bibliografia}{
	colback=gray!5!white,
	colframe=gray!75!black,
	title=Bibliografías útiles,
	fonttitle=\bfseries,
	breakable
}

% Socrative para óptica 2

\newtcolorbox{Socrative}{
	colback=orange!5!white,
	colframe=orange!75!black,
	title=Socrative,
	fonttitle=\bfseries,
	breakable
}

% Float

\usepackage{float}

% --- Hipervínculos ---
\usepackage{hyperref}
\hypersetup{
	colorlinks=true,
	linkcolor=blue,
	filecolor=magenta,      
	urlcolor=cyan,
}

% --- Datos del Documento ---
\title{Stadistical Physics}
\author{Chengyu Jin}
\date{\today}

\begin{document}
	
	\maketitle
	
	\tableofcontents
	\newpage
	
	\chapter{Solids}
	
	\section{Solids with N atoms in 3D}
	
	$T = \frac{1}{2} m \sum_{i=1}^{3N} \dot{x}_i^2 \quad ; \quad x_i \text{ : positions of } N \text{ atoms at time } t$
	
	\esp
	
	$\hookrightarrow T = \frac{1}{2} m \sum_{i=1}^{3N} \dot{u}_i^2 \quad \text{taking } u_i(t) = (x_i - x_i^0)$
	
	\esp
	
	$V(\{x_i\}) = V(\{x_i^0\}) + \sum_{i=1}^{3N} \left( \frac{\partial V}{\partial x_i} \right)_{\{x_i^0\}} (x_i - x_i^0) + \frac{1}{2} \sum_{i,j=1}^{3N} \left( \frac{\partial^2 V}{\partial x_i \partial x_j} \right)_{\{x_i^0\}} (x_i - x_i^0)(x_j - x_j^0) + \dots$
	
	\esp
	
	At the \textit{equilibrium position}: $V(\{x_i^0\}) = 0$ and $\left( \frac{\partial V}{\partial x_i} \right)_{\{x_i^0\}} = 0$.
	For small oscillations, higher order terms are neglected.
	
	\esp
	
	\begin{equation}
		H = \sum_{i=1}^{3N} \frac{1}{2} m \dot{u}_i^2 + \frac{1}{2} \sum_{i,j=1}^{3N} K_{ij} u_i u_j
	\end{equation}
	
	where $K_{ij} = \left( \frac{\partial^2 V}{\partial x_i \partial x_j} \right)_{\{x_i^0\}} \rightarrow$ Microscopic physics of the solid (its interactions).
	
	\esp
	
	But $H \neq \sum H_i$ due to the potential. As our framework needs $H = \sum H_i$, we perform a transformation to decouple it.
	

	\subsection{Normal mode decomposition}
	
	$u = \begin{pmatrix} u_1 \\ \vdots \\ u_{3N} \end{pmatrix}$. Then $T = \frac{1}{2} m \dot{u}^T \dot{u}$ and $V = \frac{1}{2} u^T K u$ (Our original Hamiltonian).
	
	\esp
	
	We search for an orthogonal transformation ($S^T S = I$) over $u = S q$ such that $S^T K S = \Lambda$, where $\Lambda$ is diagonal.
	
	\esp
	
	\begin{align*}
		\dot{u}^T &= \dot{q}^T S^T \\
		u &= S q
	\end{align*}
	then $H = \frac{1}{2} m \dot{q}^T \dot{q} + \frac{1}{2} q^T \Lambda q \rightarrow H = \frac{1}{2} \sum_{i=1}^{3N} m \dot{q}_i^2 + \frac{1}{2} \sum_{i=1}^{3N} \lambda_i q_i^2$
	
	\esp
	
	Which can be expressed as $H = \sum H_i \rightarrow H$ becomes a sum of independent oscillators!!
	\begin{equation}
		H = \sum_{i=1}^{3N} \left[ \frac{p_i^2}{2m} + \frac{1}{2} m \omega_i^2 q_i^2 \right] \text{ with } p_i = m\dot{q}_i \text{ and } \omega_i = \sqrt{\frac{\lambda_i}{m}}
	\end{equation}
	$\lambda_i$ is the eigenvalue of $K$.
	
	\esp
	
	\noindent \textit{A very complicated movement can be decomposed as a sum of normal modes.}
	
	\esp
	
	\noindent \textbf{Example 1: Eq. of motion for the displacements}
	
	\esp
	
	$H = \sum_{i=1}^{3N} \frac{p_i^2}{2m} + \frac{1}{2} \sum_{i,j=1}^{3N} K_{ij} u_i u_j$ with $p_i = m\dot{u}_i$.
	
	\esp
	
	$\dot{u}_i = \frac{\partial H}{\partial p_i} = \frac{p_i}{m}$
	$\dot{p}_i = -\frac{\partial H}{\partial u_i} \rightarrow m\ddot{u}_i = -\sum_{j=1}^{3N} K_{ij} u_j \implies \ddot{u} = \frac{-K}{m} u$
	
	\esp
	
	Since $u = S q$:
	$S \ddot{q} = -S \frac{\Lambda}{m} q \implies \ddot{q} = -\frac{\Lambda}{m} q \implies \ddot{q}_i + \frac{\lambda_i}{m} q_i = 0 \implies \ddot{q}_i + \omega_i^2 q_i = 0$
	
	\esp
	
	Solution:
	$q_i(t) = a_i \cos(\omega_i t + \phi_i)$
	$u_i(t) = (S q)_i \rightarrow u_i(t) = \sum_{j=1}^{3N} S_{ij} a_j \cos(\omega_j t + \phi_j)$
	$a_j$ and $\phi_j$ are fixed by initial conditions.
	
	\esp
	
	\noindent \textbf{Example 2: 2 atoms (1D)}
	
	\esp
	
	$V(u_1, u_2) = \frac{1}{2} k u_1^2 + \frac{1}{2} c (u_1 - u_2)^2 + \frac{1}{2} k u_2^2$
	$u(t)$?
	Initial conditions: $u_1(0) = u_2(0) = \alpha$, $\dot{u}_1(0) = \dot{u}_2(0) = 0$.

	\printbibliography[title={Bibliografía}]
\end{document}